\documentclass[aspectratio=169]{beamer}
\usetheme{Madrid}
\usecolortheme{default}
\usepackage{graphicx}
\usepackage{multimedia} % For embedding mp4 videos
\usepackage{amsmath}

\title{Human Motion Animation Generation}
\subtitle{Autoregressive Context Encoding \& Spatial-Temporal Flow Matching}
\author{Your Name}
\institute{CSE 330 Machine Learning Project, BUET}
\date{\today}

\begin{document}

\maketitle

\section{Introduction}

\begin{frame}{Problem Statement}
    \textbf{Goal:} Generate highly realistic 3D human motion animations conditioned on text descriptions.
    \vspace{0.5cm}
    
    \textbf{Context \& Dataset:} 
    \begin{itemize}
        \item Utilizing the HumanML3D dataset format.
        \item Translating text prompts (e.g., "A person is walking forward") into a sequence of 22-joint articulated poses.
    \end{itemize}
\end{frame}

\section{Evolution of Model Architecture}

\begin{frame}{Previous Iterations}
    % PLACEHOLDER for your older designs
    \textbf{Early Architectures (To be expanded)}
    \begin{itemize}
        \item \textit{Details of past approach 1 (e.g., GRU-based History Encoder)}
        \item \textit{Details of past approach 2 (e.g., Latent-only state representations)}
    \end{itemize}
    
    \vspace{0.5cm}
    \begin{block}{Animation Demonstration}
        % PLACEHOLDER: Insert path to older model output mp4
        % \movie[width=6cm,height=4cm,poster]{Click to Play}{media/old_model_demo.mp4}
        [Placeholder for previous iteration MP4]
    \end{block}
\end{frame}

\section{Current Architecture}

\begin{frame}{System Overview}
    A text-conditioned autoregressive motion generator dividing generation into 4 stages per frame:
    \begin{enumerate}
        \item \textbf{Encode}: Summarize 271D motion history with text conditioning.
        \item \textbf{Predict}: Denoise next-frame state into a reduced 68D representation via Flow Matching.
        \item \textbf{Convert}: Map 68D state back into absolute global joint positions.
        \item \textbf{Accumulate}: Derive new 271D features and append to explicit position history.
    \end{enumerate}
\end{frame}

\begin{frame}{Data Representations}
    Bridging motion geometry and model efficiency:
    \vspace{0.3cm}
    
    \textbf{271D Full Frame (Encoder \& Dataset)}
    \begin{itemize}
        \item Root state (3), Root-Invariant Coords (66), 6D Rotations (132), Local Vel (66), Contacts (4).
    \end{itemize}
    
    \vspace{0.3cm}
    
    \textbf{68D Reduced Target Space (Flow Matching)}
    \begin{itemize}
        \item Model-friendly target for flow solver.
        \item Root height/velocity/yaw (5) + 21 Non-root joints RIC (63).
    \end{itemize}
\end{frame}

\begin{frame}{1. Motion History Encoder}
    \textbf{Causal Temporal Attention Network}
    \begin{itemize}
        \item \textbf{Input:} Variable-length history window $X \in \mathbb{R}^{B \times T \times 271}$.
        \item \textbf{Text Conditioning:} Pooled CLIP text embeddings modulate every timestep via shift/scale.
        \item \textbf{Mechanism:} Causal temporal self-attention with RoPE + gated MLP blocks.
        \item \textbf{Output:} 22 joint tokens supplying context without forcing the predictor to model full spatiotemporal attention.
    \end{itemize}
\end{frame}

\begin{frame}{2. Flow Matching Predictor}
    \textbf{Spatial Transformer Engine}
    \begin{itemize}
        \item \textbf{Inputs (per token):} Noisy 68D features + History Context + 257D causal previous-frame features.
        \item \textbf{Conditioning:} AdaLN (Time $t$ embedding + CLIP Text projection).
        \item \textbf{Structure Prior:} Kinematic chain embeddings define skeletal hierarchy (depth + chain ID).
        \item \textbf{Action:} Predicts velocity fields to regress the denoised 68D state.
    \end{itemize}
\end{frame}

\begin{frame}{Training \& Inference}
    \begin{columns}
        \column{0.5\textwidth}
        \textbf{Training Loop}
        \begin{itemize}
            \item Incremental flow loss optimization.
            \item Parallel sequences using $X_{:-1}$ to predict $X_{1:}$.
            \item Uses interpolated state $x_t = t \cdot x_1 + (1-t) \cdot x_0$.
        \end{itemize}
        
        \column{0.5\textwidth}
        \textbf{Inference (Position-First)}
        \begin{itemize}
            \item Integrated using end-biased Heun ODE solver.
            \item Absolute positions act as the primary rolling state.
            \item Option to utilize Forward Kinematics (FK) consistency for skeletal stability.
        \end{itemize}
    \end{columns}
\end{frame}

\section{Results \& Demonstrations}

\begin{frame}{Generation Comparisons}
    % \movie takes [options]{poster text}{video_file}
    
    \begin{columns}[T] % align columns
        \begin{column}{.48\textwidth}
            \textbf{Ground Truth} \\
            % \movie[width=\textwidth,height=3.5cm,poster]{Play GT}{media/gt_demo.mp4}
            [Placeholder: gt\_demo.mp4]
        \end{column}
        \begin{column}{.48\textwidth}
            \textbf{Full Autoregressive Rollout} \\
            % \movie[width=\textwidth,height=3.5cm,poster]{Play Rollout}{media/rollout_demo.mp4}
            [Placeholder: rollout\_demo.mp4]
        \end{column}
    \end{columns}
\end{frame}

\begin{frame}{Teacher-Forced Comparisons}
    \begin{columns}[T]
        \begin{column}{.48\textwidth}
            \textbf{1-Step Teacher Force} \\
            % \movie[width=\textwidth,height=3.5cm,poster]{Play 1-Step TF}{media/tf_1step.mp4}
            [Placeholder: tf\_1step.mp4]
        \end{column}
        \begin{column}{.48\textwidth}
            \textbf{Masked Teacher Force} \\
            % \movie[width=\textwidth,height=3.5cm,poster]{Play Masked TF}{media/tf_masked.mp4}
            [Placeholder: tf\_masked.mp4]
        \end{column}
    \end{columns}
\end{frame}

\section{Conclusion \& Future Work}

\begin{frame}{Conclusion}
    \begin{itemize}
        \item Successfully separated causal temporal accumulation and spatial-temporal next-frame denoising.
        \item Retained geometric fidelity by using a position-first recursive inference state.
        \item Future Work: Enable scheduled sequence rollouts during training and cache state optimizations.
    \end{itemize}
    \vspace{0.8cm}
    \centering \Large \textbf{Questions?}
\end{frame}

\end{document}